\chapter[Adjuncions]{Adjuncions}
\label{cap:adjuncions}

\begin{objectius}
\apartat{Prerequisits} Functors i representabilitat (cap.~\ref{cap:functors},~\ref{cap:yoneda}); límits i colímits (cap.~\ref{cap:limits}).
\apartat{Objectius} En acabar el capítol, el lector ha de saber: definir una adjunció mitjançant la bijecció d'homconjunts i mitjançant unitat i counitat; verificar les identitats triangulars; demostrar que els adjunts són únics llevat d'isomorfisme; i aplicar que els adjunts per la dreta preserven límits i els adjunts per l'esquerra, colímits.
\end{objectius}

\section{Cap a la noció d'adjunció}

Al llarg del llibre hem trobat parelles de functors que apareixen sempre juntes: el functor lliure $\Free\colon\cat{Graph}\to\cat{Cat}$ i el functor oblit $U\colon\cat{Cat}\to\cat{Graph}$; el functor que oblida l'estructura d'un grup i el que construeix el grup lliure sobre un conjunt. En cada cas, un dels functors \enquote{afegeix estructura de la manera més econòmica possible} i l'altre \enquote{l'oblida}, i tots dos estan lligats per una correspondència precisa. Aquesta relació s'anomena \textbf{adjunció}, i és un dels conceptes centrals de la teoria de categories.

La idea es pot copsar primer en el cas més simple, el dels preordres, on una adjunció és una parella d'aplicacions monòtones lligades per una equivalència de desigualtats. Després la generalitzarem a categories qualssevol, substituint les desigualtats per bijeccions entre conjunts de morfismes.

\section{Adjuncions entre preordres}

Siguin $(P,\leq)$ i $(Q,\leq)$ dos preordres, vistos com a categories, i $f\colon P\to Q$, $g\colon Q\to P$ dues aplicacions monòtones.

\begin{definicio}\index{adjunció!entre preordres}
Diem que $f$ és \textbf{adjunta per l'esquerra} de $g$ (i $g$ \textbf{adjunta per la dreta} de $f$), i escrivim $f\dashv g$, quan per a tot $x\in P$ i tot $y\in Q$,
\[
f(x)\leq y
\quad\Longleftrightarrow\quad
x\leq g(y).
\]
\end{definicio}

\begin{exemple}[Conjunció i implicació]\index{adjunció!conjunció i implicació}
Sigui $P$ l'àlgebra de proposicions amb la relació de deduïbilitat, i fixem una fórmula $A$. Considerem les operacions
\[
f(B)=A\wedge B,\qquad g(B)=A\to B.
\]
Totes dues són monòtones. A més, $f\dashv g$, perquè
\[
A\wedge B\vdash C
\quad\Longleftrightarrow\quad
B\vdash A\to C.
\]
Aquesta equivalència és exactament la regla d'introducció i eliminació de la implicació: \emph{la conjunció amb $A$ és adjunta per l'esquerra de la implicació des de $A$}. És el primer indici que els connectius lògics són adjunts, idea que dominarà la lògica categòrica.
\end{exemple}

\begin{exemple}[Interior topològic]\index{adjunció!interior topològic}
Sigui $X$ un espai topològic, $\mathcal P(X)$ el preordre de les parts per inclusió i $\mathcal O(X)$ el dels oberts. La inclusió $i\colon\mathcal O(X)\to\mathcal P(X)$ té com a adjunt per la dreta l'\textbf{interior} $\mathrm{int}\colon\mathcal P(X)\to\mathcal O(X)$:
\[
i(U)\subseteq S
\quad\Longleftrightarrow\quad
U\subseteq\mathrm{int}(S),
\]
per a tot obert $U$ i tota part $S$. És a dir, l'interior és el més gran obert contingut en $S$.
\end{exemple}

\section{Adjuncions entre categories}

La definició general substitueix l'equivalència de desigualtats per una \textbf{bijecció entre conjunts de morfismes}, natural en totes dues variables.

\begin{definicio}[Adjunció]\index{adjunció}
\label{def:adjuncio}
Siguin $F\colon\cat{C}\to\cat{D}$ i $G\colon\cat{D}\to\cat{C}$ dos functors. Una \textbf{adjunció} $F\dashv G$ és una bijecció
\[
\theta_{A,B}\colon\Hom_{\cat{D}}(F(A),B)\;\xrightarrow{\ \cong\ }\;\Hom_{\cat{C}}(A,G(B)),
\]
per a cada $A$ de $\cat{C}$ i $B$ de $\cat{D}$, \textbf{natural} en $A$ i en $B$. Es diu que $F$ és l'\textbf{adjunt per l'esquerra} i $G$ l'\textbf{adjunt per la dreta}.
\end{definicio}

\begin{observacio}
La naturalitat vol dir que $\theta$ és un isomorfisme natural entre els dos functors
\[
\Hom_{\cat{D}}(F(-),-)\quad\text{i}\quad\Hom_{\cat{C}}(-,G(-)),
\]
tots dos de $\cat{C}\op\times\cat{D}$ cap a $\cat{Set}$. Explícitament, per a $a\colon A'\to A$ a $\cat{C}$ i $b\colon B\to B'$ a $\cat{D}$, el pas per $\theta$ commuta amb la precomposició/postcomposició per $F(a),b$ d'una banda i $a,G(b)$ de l'altra. Els morfismes que es corresponen sota $\theta$ s'anomenen \textbf{transposats} l'un de l'altre: si $\varphi\colon F(A)\to B$, el seu transposat és $\theta(\varphi)\colon A\to G(B)$.
\end{observacio}

\begin{observacio}
El cas dels preordres n'és un cas particular: si $P,Q$ són preordres, un conjunt de morfismes $\Hom(x,y)$ té com a molt un element, de manera que la bijecció $\Hom(f(x),y)\cong\Hom(x,g(y))$ es redueix a l'equivalència $f(x)\leq y\Leftrightarrow x\leq g(y)$. Així, la definició general estén literalment la dels preordres.
\end{observacio}

\section{Unitat i counitat}

Una adjunció es pot descriure, de manera equivalent, sense esmentar la bijecció, mitjançant dues transformacions naturals distingides.

\begin{definicio}[Unitat i counitat]\index{adjunció!unitat}\index{adjunció!counitat}
\label{def:unitat-counitat}
Donada una adjunció $F\dashv G$ amb bijecció $\theta$:
\begin{itemize}
\item la \textbf{unitat} és la transformació natural $\eta\colon\id_{\cat{C}}\Rightarrow G\comp F$ de components
\[
\eta_A=\theta_{A,F(A)}(\id_{F(A)})\colon A\to G(F(A));
\]
\item la \textbf{counitat} és la transformació natural $\varepsilon\colon F\comp G\Rightarrow\id_{\cat{D}}$ de components
\[
\varepsilon_B=\theta_{G(B),B}^{-1}(\id_{G(B)})\colon F(G(B))\to B.
\]
\end{itemize}
\end{definicio}

\begin{observacio}
La unitat i la counitat permeten recuperar la bijecció. Donat $\varphi\colon F(A)\to B$, el seu transposat és
\[
\theta(\varphi)=G(\varphi)\comp\eta_A\colon A\to G(B),
\]
i dualment, donat $\psi\colon A\to G(B)$, el seu transposat invers és
\[
\theta^{-1}(\psi)=\varepsilon_B\comp F(\psi)\colon F(A)\to B.
\]
La unitat expressa una \textbf{propietat universal}: per a cada $A$, el morfisme $\eta_A\colon A\to G(F(A))$ és universal des de $A$ cap a $G$, en el sentit que tot $\psi\colon A\to G(B)$ factoritza de manera única com $G(\varphi)\comp\eta_A$. Aquesta és la formulació de l'adjunció que connecta directament amb els objectes lliures.
\end{observacio}

\begin{exemple}[El monoide lliure]\index{adjunció!lliure-oblit}
Sigui $U\colon\cat{Mon}\to\cat{Set}$ el functor oblit dels monoides i $F\colon\cat{Set}\to\cat{Mon}$ el functor que envia un conjunt $X$ al \textbf{monoide lliure} $X^*$ de les paraules finites en $X$, amb la concatenació. Aleshores $F\dashv U$: un homomorfisme de monoides $X^*\to M$ està determinat de manera única pels valors sobre les lletres, és a dir, per una aplicació $X\to U(M)$. Aquesta és exactament la bijecció
\[
\Hom_{\cat{Mon}}(X^*,M)\cong\Hom_{\cat{Set}}(X,U(M)).
\]
La unitat $\eta_X\colon X\to U(X^*)$ envia cada element a la paraula d'una lletra; és la propietat universal del monoide lliure. Aquest patró ---\emph{lliure $\dashv$ oblit}--- es repeteix per a grups, espais vectorials, categories i moltes altres estructures.
\end{exemple}

\section{Identitats triangulars}

La unitat i la counitat no són independents: satisfan dues equacions que, de fet, caracteritzen les adjuncions.

\begin{proposicio}[Identitats triangulars]\index{adjunció!identitats triangulars}
Per a una adjunció $F\dashv G$ amb unitat $\eta$ i counitat $\varepsilon$, es compleixen les \textbf{identitats triangulars}:
\[
\varepsilon_{F(A)}\comp F(\eta_A)=\id_{F(A)},
\qquad
G(\varepsilon_B)\comp\eta_{G(B)}=\id_{G(B)}.
\]
Recíprocament, dues transformacions naturals $\eta\colon\id\Rightarrow G F$ i $\varepsilon\colon F G\Rightarrow\id$ que les satisfacin defineixen una adjunció $F\dashv G$.
\end{proposicio}

\begin{prova}
Comprovem la primera identitat. Pel càlcul del transposat, $\theta^{-1}(\psi)=\varepsilon_B\comp F(\psi)$; prenent $B=F(A)$ i $\psi=\eta_A=\theta(\id_{F(A)})$,
\[
\varepsilon_{F(A)}\comp F(\eta_A)=\theta^{-1}(\eta_A)=\theta^{-1}(\theta(\id_{F(A)}))=\id_{F(A)}.
\]
La segona identitat és dual, usant $\theta(\varphi)=G(\varphi)\comp\eta_A$ amb $A=G(B)$ i $\varphi=\varepsilon_B$. Per al recíproc, es defineix la bijecció per $\theta(\varphi)=G(\varphi)\comp\eta_A$ i $\theta^{-1}(\psi)=\varepsilon_B\comp F(\psi)$; les identitats triangulars garanteixen que són inverses l'una de l'altra, i la naturalitat de $\eta,\varepsilon$ dona la de $\theta$.
\end{prova}

\begin{observacio}
Les identitats triangulars reben aquest nom perquè expressen la commutativitat de dos triangles que involucren $F$, $G$, $\eta$ i $\varepsilon$. Donen una definició d'adjunció \emph{purament diagramàtica}, sense referència a conjunts de morfismes, la qual cosa és sovint la més còmoda per treballar i l'única disponible en contextos on no hi ha homconjunts (per exemple, en 2-categories abstractes).
\end{observacio}

\section{Unicitat dels adjunts}

\begin{proposicio}\index{adjunció!unicitat}
Els adjunts són únics llevat d'isomorfisme natural. Si $F\dashv G$ i $F\dashv G'$, aleshores $G\cong G'$; i si $F\dashv G$ i $F'\dashv G$, aleshores $F\cong F'$.
\end{proposicio}

\begin{prova}
Si $F\dashv G$ i $F\dashv G'$, aleshores per a tot $A,B$,
\[
\Hom(A,G(B))\cong\Hom(F(A),B)\cong\Hom(A,G'(B)),
\]
naturalment en $A$. Pel corol·lari del lema de Yoneda ---dos objectes amb functors representables naturalment isomorfs són isomorfs---, $G(B)\cong G'(B)$ naturalment en $B$, és a dir, $G\cong G'$. El cas de $F$ és dual.
\end{prova}

\section{Els adjunts i els límits}

El resultat més útil sobre adjuncions diu com interactuen amb límits i colímits. És la raó de fons per la qual els functors oblit preserven productes però no coproductes.

\begin{teorema}[Els adjunts per la dreta preserven límits]\index{adjunció!preservació de límits}
\label{teo:adjunts-preserven-limits}
Si $G\colon\cat{D}\to\cat{C}$ és un adjunt per la dreta, aleshores $G$ preserva tots els límits que existeixen a $\cat{D}$. Dualment, un adjunt per l'esquerra preserva tots els colímits.
\end{teorema}

\begin{prova}[Idea]
Sigui $F\dashv G$ i sigui $(L,\lambda)$ un límit d'un diagrama $D\colon\cat{I}\to\cat{D}$. Per a tot objecte $A$ de $\cat{C}$, usant l'adjunció i que $\Hom(A,-)$ transforma límits,
\[
\Hom(A,G(L))\cong\Hom(F(A),L)\cong\varprojlim\Hom(F(A),D_i)\cong\varprojlim\Hom(A,G(D_i)).
\]
Això diu que $G(L)$ representa el functor dels cons sobre $G\comp D$, és a dir, que $(G(L),G\lambda)$ n'és el límit. Per tant $G$ preserva el límit. El cas dels colímits és dual.
\end{prova}

\begin{observacio}
Aquest teorema es recorda amb les sigles angleses \enquote{RAPL} (\emph{right adjoints preserve limits}) i \enquote{LAPC} (\emph{left adjoints preserve colimits}). Explica de cop molts fets: els functors oblit, que solen ser adjunts per la dreta, preserven productes, equalitzadors i límits en general, però no coproductes; i els functors lliures, adjunts per l'esquerra, preserven coproductes però no productes. És també un criteri negatiu útil: un functor que \emph{no} preserva límits no pot ser un adjunt per la dreta.
\end{observacio}

\subsection{El límit i el colímit com a adjunts del functor diagonal}\index{límit!com a adjunt}\index{colímit!com a adjunt}

Hi ha una segona manera, més sintètica, en què límits i colímits són adjuncions ---no ja \emph{preservats} per adjunts, sinó ells mateixos adjunts. Fixem una categoria d'índexs $\cat{I}$ i considerem el \textbf{functor diagonal}
\[
\Delta\colon\cat{C}\longrightarrow\cat{C}^{\cat{I}},
\]
que envia un objecte $A$ al diagrama constant de valor $A$ (i un morfisme a la transformació natural constant). Prendre el límit i prendre el colímit d'un diagrama són operacions $\cat{C}^{\cat{I}}\to\cat{C}$, i resulten ser, respectivament, l'adjunt per la dreta i per l'esquerra de $\Delta$.

\begin{teorema}[$\colimop\dashv\Delta\dashv\varprojlim$]\index{adjunció!límit-diagonal}
\label{teo:lim-adjunt}
Sigui $\cat{I}$ una categoria petita i suposem que $\cat{C}$ té tots els límits i colímits de forma $\cat{I}$. Aleshores el functor diagonal $\Delta\colon\cat{C}\to\cat{C}^{\cat{I}}$ té adjunt per l'esquerra $\colimop_{\cat{I}}$ i adjunt per la dreta $\varprojlim_{\cat{I}}$:
\[
\colimop_{\cat{I}}\;\dashv\;\Delta\;\dashv\;\varprojlim_{\cat{I}}.
\]
\end{teorema}

\begin{prova}[Idea]
Desplegem l'adjunció de la dreta. Un morfisme $\Delta A\Rightarrow D$ a $\cat{C}^{\cat{I}}$ és, per definició, una família compatible $A\to D_i$, és a dir, un con sobre $D$ amb vèrtex $A$; i la propietat universal del límit diu exactament que aquests cons estan en bijecció natural amb els morfismes $A\to\varprojlim D$:
\[
\Hom_{\cat{C}^{\cat{I}}}(\Delta A, D)\;\cong\;\Hom_{\cat{C}}(A,\varprojlim D),
\]
naturalment en $A$ i en $D$. Això és precisament $\Delta\dashv\varprojlim$. Dualment, un morfisme $D\Rightarrow\Delta A$ és un cocon, i la propietat universal del colímit dona $\Hom_{\cat{C}}(\colimop D,A)\cong\Hom_{\cat{C}^{\cat{I}}}(D,\Delta A)$, és a dir, $\colimop\dashv\Delta$.
\end{prova}

\begin{observacio}
Aquesta fórmula és una de les síntesis més netes del llibre: recull en una sola línia tot el capítol de límits. La unitat de $\Delta\dashv\varprojlim$ és la família de projeccions del límit; la counitat de $\colimop\dashv\Delta$, la d'injeccions del colímit. I la lectura és profètica: quan al capítol de subobjectes vegem l'existencial i l'universal com a adjunts de la reindexació $f^{*}$, estarem veient una versió \enquote{fibrada} d'aquesta mateixa idea ---$f^{*}$ fa el paper de $\Delta$, i els quantificadors, el de $\colimop$ i $\varprojlim$.
\end{observacio}

\section{Exponencials i currificació}

Recuperem ara la relació que vam ajornar en tractar les categories de functors: la que lliga el producte amb els objectes de functors.

\begin{proposicio}\index{adjunció!producte-exponencial}\index{Cat@$\cat{Cat}$!exponencials}
Per a categories petites $\cat{A},\cat{B},\cat{C}$, hi ha un isomorfisme natural
\[
\Hom_{\cat{Cat}}(\cat{A}\times\cat{B},\cat{C})\;\cong\;\Hom_{\cat{Cat}}(\cat{A},\cat{C}^{\cat{B}}).
\]
És a dir, el functor $-\times\cat{B}$ és adjunt per l'esquerra del functor $(-)^{\cat{B}}$:
\[
-\times\cat{B}\;\dashv\;(-)^{\cat{B}}.
\]
\end{proposicio}

\begin{prova}[Idea]
Un functor $\cat{A}\times\cat{B}\to\cat{C}$ es pot \enquote{currificar}: fixant la primera variable en cada objecte $A$ de $\cat{A}$, s'obté un functor $\cat{B}\to\cat{C}$, i aquesta assignació és functorial en $A$, de manera que dona un functor $\cat{A}\to\cat{C}^{\cat{B}}$. La correspondència inversa \enquote{descurrifica}. Que sigui una bijecció natural es comprova component a component, i és el reflex, a nivell de categories, de la bijecció conjuntista $\Hom(A\times B,C)\cong\Hom(A,C^B)$.
\end{prova}

\begin{observacio}
Aquesta adjunció ---\emph{producte $\dashv$ exponencial}--- és el prototip de les \textbf{categories cartesianes tancades}, que estudiarem al capítol següent i que són el marc categòric del càlcul lambda simplement tipat. La \enquote{currificació} és exactament la que apareix en programació funcional: transformar una funció de dues variables en una funció que retorna funcions.
\end{observacio}

\section{Cap al tancament cartesià i la lògica}

Les adjuncions han unificat construccions aparentment diverses: objectes lliures, interiors i clausures, connectius lògics, exponencials. El patró comú ---una bijecció natural d'homconjunts, o equivalentment una unitat i una counitat amb les identitats triangulars--- reapareixerà com a columna vertebral de tot el que segueix.

Al capítol següent, dedicat a les \textbf{categories cartesianes tancades}, l'adjunció entre el producte i l'exponencial es fa interna a la categoria i dona lloc als objectes exponencials i a la interpretació categòrica del càlcul lambda. I més endavant, en entrar de ple a la lògica categòrica, veurem que els \textbf{quantificadors} $\exists$ i $\forall$ són, respectivament, els adjunts per l'esquerra i per la dreta dels functors de reindexació entre subobjectes ---la formulació categòrica precisa de la relació entre quantificació i substitució. Les adjuncions són, doncs, el pont directe cap a la semàntica categòrica de la lògica.

\begin{resumcap}
\apartat{Cinc idees} (1)~Una adjunció $F\dashv G$ és una bijecció natural $\Hom(F A,B)\cong\Hom(A,G B)$. (2)~Equivale a donar unitat i counitat que satisfan les identitats triangulars. (3)~La unitat expressa una propietat universal (morfisme universal cap a $G$). (4)~Els adjunts són únics llevat d'isomorfisme. (5)~Els adjunts per la dreta preserven límits; els adjunts per l'esquerra, colímits.
\apartat{Errors típics} Intercanviar quin adjunt preserva límits i quin colímits; oblidar la naturalitat de la bijecció d'homconjunts; confondre unitat i counitat; i tractar les mnemotècniques (\enquote{RAPL\,/\,LAPC}) com a terminologia establerta en lloc d'ajudes informals.
\apartat{Lectura recomanada} \cite[cap.~9]{awodey}; \cite[cap.~2]{leinster} per a una exposició centrada en adjuncions; \cite[IV]{maclane} per al tractament clàssic i \cite{borceux} per a més detalls.
\end{resumcap}
